\documentclass[11pt]{article}
\newcommand{\ListaTitulo}{Gabarito --- Lista 08}
% Preâmbulo compartilhado — MATD48, Listas de Exercícios 2026
% Usado por Lista{NN}.tex e Gabarito{NN}.tex via % Preâmbulo compartilhado — MATD48, Listas de Exercícios 2026
% Usado por Lista{NN}.tex e Gabarito{NN}.tex via % Preâmbulo compartilhado — MATD48, Listas de Exercícios 2026
% Usado por Lista{NN}.tex e Gabarito{NN}.tex via \input{preamble}

\usepackage[utf8]{inputenc}
\usepackage[T1]{fontenc}
\usepackage[brazil]{babel}
\usepackage{fullpage}
\usepackage{amsmath,amssymb}
\usepackage{enumitem}
\usepackage{xcolor}
\usepackage{fancyhdr}
\usepackage{titlesec}
\usepackage{listings}
\usepackage{hyperref}
\usepackage{booktabs}
\usepackage{graphicx}

\definecolor{matd48azul}{HTML}{0B4F6C}
\definecolor{matd48verde}{HTML}{2E8B57}

\titleformat{\section}{\normalfont\Large\bfseries\color{matd48azul}}{\thesection}{1em}{}
\titleformat{\subsection}{\normalfont\large\bfseries\color{matd48azul}}{\thesubsection}{1em}{}

\providecommand{\ListaTitulo}{Lista de Exercícios}

\setlength{\headheight}{14pt}
\pagestyle{fancy}
\fancyhf{}
\lhead{\small MATD48 --- Planejamento de Experimentos A}
\rhead{\small \ListaTitulo}
\cfoot{\thepage}
\renewcommand{\headrulewidth}{0.4pt}

\lstset{
  basicstyle=\ttfamily\small,
  breaklines=true,
  frame=single,
  columns=fullflexible,
  backgroundcolor=\color{gray!8},
  commentstyle=\color{matd48verde},
  keywordstyle=\color{matd48azul}\bfseries,
  language=R,
  extendedchars=true,
  literate=%
    {á}{{\'a}}1 {à}{{\`a}}1 {â}{{\^a}}1 {ã}{{\~a}}1
    {é}{{\'e}}1 {ê}{{\^e}}1 {í}{{\'i}}1 {ó}{{\'o}}1
    {ô}{{\^o}}1 {õ}{{\~o}}1 {ú}{{\'u}}1 {ç}{{\c c}}1
    {Á}{{\'A}}1 {À}{{\`A}}1 {Â}{{\^A}}1 {Ã}{{\~A}}1
    {É}{{\'E}}1 {Ê}{{\^E}}1 {Í}{{\'I}}1 {Ó}{{\'O}}1
    {Ô}{{\^O}}1 {Õ}{{\~O}}1 {Ú}{{\'U}}1 {Ç}{{\c C}}1
}

\hypersetup{colorlinks=true, linkcolor=matd48azul, urlcolor=matd48azul, citecolor=matd48azul}

\newcommand{\cabecalho}[2]{%
  \begin{center}
    {\Large\bfseries\color{matd48azul} #1}\\[0.3em]
    {\large #2}\\[0.3em]
    \rule{\linewidth}{0.6pt}
  \end{center}
  \vspace{0.5em}
}

\newenvironment{questao}[1]{\vspace{1em}\noindent\textbf{Questão #1.}\hspace{0.4em}}{\par}
\newenvironment{solucao}{\vspace{0.3em}\noindent\textit{Solução.}\hspace{0.4em}\color{black!85}}{\par\vspace{0.5em}}


\usepackage[utf8]{inputenc}
\usepackage[T1]{fontenc}
\usepackage[brazil]{babel}
\usepackage{fullpage}
\usepackage{amsmath,amssymb}
\usepackage{enumitem}
\usepackage{xcolor}
\usepackage{fancyhdr}
\usepackage{titlesec}
\usepackage{listings}
\usepackage{hyperref}
\usepackage{booktabs}
\usepackage{graphicx}

\definecolor{matd48azul}{HTML}{0B4F6C}
\definecolor{matd48verde}{HTML}{2E8B57}

\titleformat{\section}{\normalfont\Large\bfseries\color{matd48azul}}{\thesection}{1em}{}
\titleformat{\subsection}{\normalfont\large\bfseries\color{matd48azul}}{\thesubsection}{1em}{}

\providecommand{\ListaTitulo}{Lista de Exercícios}

\setlength{\headheight}{14pt}
\pagestyle{fancy}
\fancyhf{}
\lhead{\small MATD48 --- Planejamento de Experimentos A}
\rhead{\small \ListaTitulo}
\cfoot{\thepage}
\renewcommand{\headrulewidth}{0.4pt}

\lstset{
  basicstyle=\ttfamily\small,
  breaklines=true,
  frame=single,
  columns=fullflexible,
  backgroundcolor=\color{gray!8},
  commentstyle=\color{matd48verde},
  keywordstyle=\color{matd48azul}\bfseries,
  language=R,
  extendedchars=true,
  literate=%
    {á}{{\'a}}1 {à}{{\`a}}1 {â}{{\^a}}1 {ã}{{\~a}}1
    {é}{{\'e}}1 {ê}{{\^e}}1 {í}{{\'i}}1 {ó}{{\'o}}1
    {ô}{{\^o}}1 {õ}{{\~o}}1 {ú}{{\'u}}1 {ç}{{\c c}}1
    {Á}{{\'A}}1 {À}{{\`A}}1 {Â}{{\^A}}1 {Ã}{{\~A}}1
    {É}{{\'E}}1 {Ê}{{\^E}}1 {Í}{{\'I}}1 {Ó}{{\'O}}1
    {Ô}{{\^O}}1 {Õ}{{\~O}}1 {Ú}{{\'U}}1 {Ç}{{\c C}}1
}

\hypersetup{colorlinks=true, linkcolor=matd48azul, urlcolor=matd48azul, citecolor=matd48azul}

\newcommand{\cabecalho}[2]{%
  \begin{center}
    {\Large\bfseries\color{matd48azul} #1}\\[0.3em]
    {\large #2}\\[0.3em]
    \rule{\linewidth}{0.6pt}
  \end{center}
  \vspace{0.5em}
}

\newenvironment{questao}[1]{\vspace{1em}\noindent\textbf{Questão #1.}\hspace{0.4em}}{\par}
\newenvironment{solucao}{\vspace{0.3em}\noindent\textit{Solução.}\hspace{0.4em}\color{black!85}}{\par\vspace{0.5em}}


\usepackage[utf8]{inputenc}
\usepackage[T1]{fontenc}
\usepackage[brazil]{babel}
\usepackage{fullpage}
\usepackage{amsmath,amssymb}
\usepackage{enumitem}
\usepackage{xcolor}
\usepackage{fancyhdr}
\usepackage{titlesec}
\usepackage{listings}
\usepackage{hyperref}
\usepackage{booktabs}
\usepackage{graphicx}

\definecolor{matd48azul}{HTML}{0B4F6C}
\definecolor{matd48verde}{HTML}{2E8B57}

\titleformat{\section}{\normalfont\Large\bfseries\color{matd48azul}}{\thesection}{1em}{}
\titleformat{\subsection}{\normalfont\large\bfseries\color{matd48azul}}{\thesubsection}{1em}{}

\providecommand{\ListaTitulo}{Lista de Exercícios}

\setlength{\headheight}{14pt}
\pagestyle{fancy}
\fancyhf{}
\lhead{\small MATD48 --- Planejamento de Experimentos A}
\rhead{\small \ListaTitulo}
\cfoot{\thepage}
\renewcommand{\headrulewidth}{0.4pt}

\lstset{
  basicstyle=\ttfamily\small,
  breaklines=true,
  frame=single,
  columns=fullflexible,
  backgroundcolor=\color{gray!8},
  commentstyle=\color{matd48verde},
  keywordstyle=\color{matd48azul}\bfseries,
  language=R,
  extendedchars=true,
  literate=%
    {á}{{\'a}}1 {à}{{\`a}}1 {â}{{\^a}}1 {ã}{{\~a}}1
    {é}{{\'e}}1 {ê}{{\^e}}1 {í}{{\'i}}1 {ó}{{\'o}}1
    {ô}{{\^o}}1 {õ}{{\~o}}1 {ú}{{\'u}}1 {ç}{{\c c}}1
    {Á}{{\'A}}1 {À}{{\`A}}1 {Â}{{\^A}}1 {Ã}{{\~A}}1
    {É}{{\'E}}1 {Ê}{{\^E}}1 {Í}{{\'I}}1 {Ó}{{\'O}}1
    {Ô}{{\^O}}1 {Õ}{{\~O}}1 {Ú}{{\'U}}1 {Ç}{{\c C}}1
}

\hypersetup{colorlinks=true, linkcolor=matd48azul, urlcolor=matd48azul, citecolor=matd48azul}

\newcommand{\cabecalho}[2]{%
  \begin{center}
    {\Large\bfseries\color{matd48azul} #1}\\[0.3em]
    {\large #2}\\[0.3em]
    \rule{\linewidth}{0.6pt}
  \end{center}
  \vspace{0.5em}
}

\newenvironment{questao}[1]{\vspace{1em}\noindent\textbf{Questão #1.}\hspace{0.4em}}{\par}
\newenvironment{solucao}{\vspace{0.3em}\noindent\textit{Solução.}\hspace{0.4em}\color{black!85}}{\par\vspace{0.5em}}


\begin{document}

\cabecalho{Gabarito --- Lista de Exercícios 08}{Análise de covariância e teste de Kruskal-Wallis (Capítulo 3, segunda metade / Aula 08)}

\begin{questao}{1} \textbf{(Psicologia --- médias ajustadas em ANCOVA)}
\end{questao}
\begin{solucao}
\begin{enumerate}[label=(\alph*)]
  \item $\bar y_{0(\text{aj})} = 91{,}0 - 0{,}60(57-60) = 91{,}0 - 0{,}60(-3) = 91{,}0+1{,}8 =
    \mathbf{92{,}8}$.
    $\bar y_{200(\text{aj})} = 108{,}0 - 0{,}60(64-60) = 108{,}0-2{,}4 = \mathbf{105{,}6}$.
  \item Diferença bruta: $108{,}0-91{,}0=17{,}0$. Diferença ajustada: $105{,}6-92{,}8=12{,}8$. O
    ajuste \textbf{diminuiu} a diferença estimada. Isso acontece porque o grupo de 200~mg tinha,
    por acaso, atenção basal \emph{acima} da média geral ($64>60$) e o grupo de 0~mg tinha basal
    \emph{abaixo} da média geral ($57<60$): parte da diferença bruta de 17 pontos não se deve à
    dose em si, mas ao fato de que o grupo de 200~mg já começava com atenção basal mais alta. A
    ANCOVA remove essa parte "emprestada" da diferença, isolando o efeito atribuível à dose.
  \item As médias ajustadas representam a comparação causal entre doses \emph{mantendo a
    covariável fixa} no seu valor médio geral, removendo o confundimento causado por
    desbalanceamento --- ainda que pequeno e devido ao acaso --- da linha de base entre grupos.
    As médias brutas confundem o efeito do tratamento com diferenças pré-existentes que a
    aleatorização não eliminou por completo nesta amostra específica (ela só garante
    comparabilidade \emph{em expectativa}, não em cada realização finita do experimento).
\end{enumerate}
\end{solucao}

\begin{questao}{2} \textbf{(Ciência de dados --- verificando a homogeneidade das inclinações)}
\end{questao}
\begin{solucao}
\begin{enumerate}[label=(\alph*)]
  \item Testa-se se a inclinação (o efeito do número de imagens sobre o tempo de carregamento) é
    a \textbf{mesma} nas quatro configurações. $H_0: \beta_1=\beta_2=\beta_3=\beta_4$ (equivalente
    a: não há interação entre a covariável e o fator configuração) --- em outras palavras, o
    quanto cada imagem adicional atrasa o carregamento não depende de qual configuração de
    scripts está ativa.
  \item Sim: $p=0{,}268 > 0{,}05$, não se rejeita $H_0$ de inclinações homogêneas. A equipe pode
    prosseguir com o modelo aditivo padrão de ANCOVA, com uma única inclinação comum $\beta$.
  \item Com $p=0{,}004$, rejeitaríamos $H_0$: as inclinações diferem entre as configurações --- o
    efeito do número de imagens sobre o tempo de carregamento depende de qual configuração está
    ativa (por exemplo, pode ser que scripts de anúncios piorem desproporcionalmente páginas já
    carregadas de imagens). Um único $\beta$ comum deixaria de representar corretamente nenhum
    grupo em particular, produzindo um ajuste tendencioso. A equipe deveria abandonar o modelo
    aditivo e, por exemplo, analisar a relação covariável-resposta separadamente por
    configuração (inclinações simples por grupo), ou reportar e interpretar a própria interação
    como um resultado substantivo do estudo.
\end{enumerate}
\end{solucao}

\begin{questao}{3} \textbf{(Agricultura --- Kruskal-Wallis, cálculo manual)}
\end{questao}
\begin{solucao}
\begin{enumerate}[label=(\alph*)]
  \item Ordenando as 9 observações: $10(1), 11(2), 12(3), 14(4), 15(5), 18(6), 20(7), 22(8),
    25(9)$. Variedade A: $12\to3,\ 18\to6,\ 15\to5$, $R_A=3+6+5=14$. Variedade B:
    $20\to7,\ 25\to9,\ 22\to8$, $R_B=7+9+8=24$. Variedade C: $10\to1,\ 14\to4,\ 11\to2$,
    $R_C=1+4+2=7$. Soma: $R_A+R_B+R_C=14+24+7=45=N(N+1)/2=9(10)/2=45$. \checkmark
  \item
    \[
    H = \frac{12}{9(10)}\left(\frac{14^2}{3}+\frac{24^2}{3}+\frac{7^2}{3}\right) - 3(10)
    = \frac{12}{90}\left(\frac{196+576+49}{3}\right) - 30
    = 0{,}1333 \times 273{,}67 - 30 \approx \mathbf{6{,}49}.
    \]
  \item $H \approx 6{,}49 > 5{,}99 = \chi^2_{0{,}05;2}$: \textbf{rejeita-se $H_0$}. Há evidência,
    ao nível de 5\%, de que as três variedades diferem quanto à distribuição (mediana) de
    produtividade.
\end{enumerate}
\end{solucao}

\begin{questao}{4} \textbf{(Dedução --- não viesamento da média ajustada em ANCOVA)}
\end{questao}
\begin{solucao}
\begin{enumerate}[label=(\alph*)]
  \item Tomando esperança em $y_{ij}=\mu+\tau_i+\beta(x_{ij}-\bar x_{\cdot\cdot})+\varepsilon_{ij}$
    (com $x_{ij}$ fixo e $E[\varepsilon_{ij}]=0$): $E[y_{ij}] = \mu+\tau_i+\beta(x_{ij}-\bar
    x_{\cdot\cdot})$. Tomando a média sobre $j=1,\ldots,n_i$ dentro do grupo $i$:
    \[
    E[\bar y_{i\cdot}] = \mu+\tau_i+\beta\left(\frac{1}{n_i}\sum_j x_{ij} - \bar
    x_{\cdot\cdot}\right) = \mu+\tau_i+\beta(\bar x_{i\cdot}-\bar x_{\cdot\cdot}) =
    \mu_i+\beta(\bar x_{i\cdot}-\bar x_{\cdot\cdot}).
    \]
  \item Como $x_{ij}$ é fixo, $\bar x_{i\cdot}-\bar x_{\cdot\cdot}$ é uma \textbf{constante}, não
    uma variável aleatória. Logo, por linearidade da esperança,
    \[
    E[\bar y_{i\cdot(\text{aj})}] = E\big[\bar y_{i\cdot} - \hat\beta(\bar x_{i\cdot}-\bar
    x_{\cdot\cdot})\big] = E[\bar y_{i\cdot}] - (\bar x_{i\cdot}-\bar x_{\cdot\cdot})\,E[\hat\beta].
    \]
    Substituindo o item (a) e $E[\hat\beta]=\beta$:
    \[
    E[\bar y_{i\cdot(\text{aj})}] = \big[\mu_i+\beta(\bar x_{i\cdot}-\bar x_{\cdot\cdot})\big] -
    (\bar x_{i\cdot}-\bar x_{\cdot\cdot})\,\beta = \mu_i. \qquad \blacksquare
    \]
  \item A demonstração usa a hipótese de que a covariável não é afetada pelo tratamento
    exatamente no passo em que tratamos $x_{ij}$ (e, portanto, $\bar x_{i\cdot}$) como
    \textbf{fixo/constante}, e não como uma variável aleatória cuja distribuição poderia depender
    de $\tau_i$. Se a covariável fosse medida \emph{depois} do tratamento (afetada por ele),
    $\bar x_{i\cdot}$ carregaria parte do próprio efeito de tratamento, e $E[\bar x_{i\cdot}]$
    dependeria de $i$ por uma razão distinta de mero acaso amostral --- nesse caso, "ajustar" por
    $x$ removeria também uma parte do efeito causal genuíno do tratamento (viés de
    superajuste, ou \emph{post-treatment bias}), e a conclusão $E[\bar y_{i\cdot(\text{aj})}] =
    \mu_i$ deixaria de valer, porque $\mu_i$ deixaria de ser a quantidade causal de interesse
    isolada corretamente pelo ajuste.
\end{enumerate}
\end{solucao}

\begin{questao}{5} \textbf{(Ciência de dados --- Kruskal-Wallis vs. ANOVA, e o efeito de empates)}
\end{questao}
\begin{solucao}
\begin{enumerate}[label=(\alph*)]
  \item A ANOVA depende de normalidade dos resíduos para a validade (ao menos aproximada) do
    teste $F$, e é sensível a valores extremos: poucas observações muito grandes (a cauda de
    conexões lentas) podem dominar a soma de quadrados e distorcer as médias e a estimativa de
    variância. O Kruskal-Wallis substitui os valores por postos, o que limita a influência de
    qualquer observação extrema a, no máximo, o maior posto disponível --- tornando o teste
    robusto à forma da distribuição e a outliers, preservando o controle do erro tipo I mesmo sem
    transformar os dados.
  \item Sem a correção por empates, $H$ tende a ficar \textbf{menor} do que deveria (o teste fica
    \emph{conservador}, não anticonservador). O fator de correção $C = 1-\sum_j(\tau_j^3-\tau_j)/
    (N^3-N)$ satisfaz $C \leq 1$, e a estatística corrigida é $H_c = H/C \geq H$. Empates
    comprimem a variabilidade possível entre postos (postos empatados recebem a mesma média de
    posto, reduzindo a dispersão), e a fórmula padrão de $H$ pressupõe implicitamente ausência de
    empates; sem dividir por $C<1$ para compensar essa compressão, o valor calculado de $H$ fica
    artificialmente pequeno relativamente à referência $\chi^2$, reduzindo o poder do teste
    (mais difícil rejeitar $H_0$ quando ela é de fato falsa) na presença de muitos empates.
  \item Duas distribuições podem ter a mesma mediana mas formas muito diferentes: por exemplo,
    um grupo com valores fortemente concentrados perto de 100~ms (baixa dispersão) e outro grupo
    também com mediana 100~ms, mas bimodal ou de cauda pesada (muitas observações bem abaixo e
    bem acima de 100~ms). Mesmo com medianas idênticas, a distribuição dos \emph{postos}
    combinados tende a diferir entre os dois grupos --- o grupo de cauda pesada contribui
    desproporcionalmente tanto para postos muito baixos quanto muito altos --- e o
    Kruskal-Wallis, que testa igualdade de distribuições (não só de posição central), pode
    detectar essa diferença de forma. Isso ilustra por que uma rejeição de $H_0$ no
    Kruskal-Wallis não deve ser automaticamente interpretada como "as medianas diferem": pode
    haver diferença de forma/dispersão em vez de (ou além de) diferença de posição central.
\end{enumerate}
\end{solucao}

\end{document}
